\chapter{Repository Workflow}

\section{Primary COMSOL Processing Workflow}

The primary processing workflow is:
\begin{enumerate}
\item export COMSOL J--V and state tables to
\texttt{data/raw/comsol/time\_series/};
\item run:
\begin{verbatim}
python .\scripts\run_comsol_paper_pipeline.py
\end{verbatim}
\item inspect processed data under \texttt{data/processed/comsol/};
\item inspect generated figures under \texttt{outputs/figures/}; and
\item rebuild the thesis with
\texttt{latexmk -pdf -interaction=nonstopmode main.tex}.
\end{enumerate}

\section{Traceability Files}

The repository includes support files for reproducibility and archival use:
\begin{itemize}
\item \texttt{README.md};
\item \texttt{START\_HERE.md};
\item \texttt{REPRODUCIBILITY.md};
\item \texttt{docs/data\_dictionary.md};
\item \texttt{docs/traceability\_matrix.md};
\item \texttt{docs/journal\_submission\_checklist.md}; and
\item \texttt{CITATION.cff}.
\end{itemize}

\section{External Calibration Context}

The external Wang et al. source-data workflow can be run as supporting
calibration context:
\begin{verbatim}
python .\scripts\run_psc_degradation_study.py
\end{verbatim}
That workflow is not the central thesis result. The central result is the
COMSOL model, the processed COMSOL scenario library, the mechanism-specific
parameter mappings, and the experimental calibration plan.

\section{Electronic Parameter-Sweep Workflow}

The uploaded COMSOL Global Evaluation sweep for $N_{t,0}$, $R_{s,0}$, and
$R_{sh,0}$ is processed separately from the aging-time pipeline:
\begin{verbatim}
python .\scripts\plot_comsol_parameter_sweep.py
\end{verbatim}
This command reads the architecture-scoped parameter-sweep export under
\path{data/raw/comsol/arch_baseline_pin/parameter_sweeps/}, extracts the active
voltage-ramp segment for each parameter combination, computes J--V metrics, and
writes sweep figures under
\path{outputs/figures/comsol_parameter_sweep/}. It also writes:
\begin{itemize}
\item \path{outputs/tables/comsol_parameter_sweep_metrics.csv}; and
\item \path{outputs/tables/comsol_parameter_sweep_single_axis_metrics.csv}.
\end{itemize}

\section{Surrogate Workbench Prototype}

The supporting surrogate application is launched with:
\begin{verbatim}
streamlit run .\apps\psc_surrogate_app\streamlit_app.py
\end{verbatim}
The current application uses placeholder prediction logic and current COMSOL
seed tables. Its purpose is to define the future data contract, run
provenance, export bundle format, inverse-diagnosis interface, and
active-learning workflow before the full COMSOL design of experiments is
available.

\section{Architecture-Stress Grid Workflow}

The architecture-scoped aging exports are processed with:
\begin{verbatim}
python .\scripts\process_comsol_architecture_stress_jv.py
python .\scripts\process_comsol_profile_grid.py
\end{verbatim}
The first command extracts diagnostic J--V curves and PV metrics from the
baseline p-i-n temperature-aging exports. The second command summarizes the
6 $\times$ 6 light-temperature profile grid into condition-time-variable rows
with spatial means, extrema, and interface-region summaries. The profile-grid
summary is the current bridge between COMSOL internal states and future
surrogate or inverse-diagnosis data sets.
